\chapter{二维运动学}\label{ch:g12:kinematics-2d}

看一辆汽车以稳定的 \qty{30}{km/h} 绕环岛行驶：指针
始终不动，但每位乘客都感到被甩向一侧——运动
在变化，变化的并不是速率。每一时刻一个数
已不够用：平面内的运动需要用箭头表示。本章
为运动的质点配备三个量——位置、速度、加速度——
它们由导数相互联系，并借助几何图形
读出整类运动的特征。

\section{位置矢量}

运动只有相对于指定的参考系才有意义——参考系由刚体
加上一只时钟构成（见\cref{ch:g10:relative-motion}）。本年度的新内容是
精确化：在参考系上取原点 $O$ 与两条互相垂直的
带刻度的坐标轴 $Ox$、$Oy$，并记录质点在每一时刻的
位置，写成一对函数。

\begin{definition}[位置矢量]\label{def:g12:kinematics-2d:position}
在配有原点 $O$ 与坐标轴 $Ox$、$Oy$ 的参考系中，
\emph{位置矢量}\index{位置矢量}即运动质点 $M$
在时刻 $t$ 的位置为 $\vect{OM}(t)$，其坐标为
$\bigl(x(t),\, y(t)\bigr)$：关于时间的两个函数，各对应一条坐标轴。
$M$ 随 $t$ 变化所描出的曲线即为它的轨迹。
\end{definition}

\section{速度矢量}

\begin{definition}[速度矢量]\label{def:g12:kinematics-2d:velocity}
在很短的时间间隔内，质点的位移为
$\vect{OM}(t+\Delta t) - \vect{OM}(t)$；除以 $\Delta t$，再令时间间隔趋于零，
即得到你在数学课程中学过的导数，
只不过分别作用于每个坐标。
\emph{速度矢量}\index{速度矢量}是 $M$ 的位置矢量的导数，
$\vect v(t) = \bigl(x'(t),\, y'(t)\bigr)$，单位为
\unit{m/s}。其模 $v = \sqrt{x'^2 + y'^2}$ 即为
\emph{速率}\index{速率}——速度计上显示的那个数。
\end{definition}

\begin{example}[无人机的水平飞行]\label{ex:g12:kinematics-2d:drone}
一架无人机的运动规律为 $x(t) = 4.0\,t$，$y(t) = 3.0\,t$（
长度单位为米，时间单位为秒）：轨迹为直线 $y = \tfrac34 x$，速度
在任意时刻均为 $\vect v = (4.0, 3.0)$，速率
为 $\sqrt{4.0^2 + 3.0^2} = \qty{5.0}{m/s}$，且保持不变。
\end{example}

\begin{proposition}[速度与轨迹相切]\label{prop:g12:kinematics-2d:tangent}
在每一时刻，$\vect v(t)$ 都与轨迹在
$M$ 当前所在位置处相切，并指向运动方向。
\end{proposition}

\begin{proof}
位移是轨迹的一条弦；当 $\Delta t$ 趋于零时，
弦的方向趋于切线的方向。
\end{proof}

\section{加速度矢量}

\begin{definition}[加速度矢量]\label{def:g12:kinematics-2d:acceleration}
\emph{加速度矢量}\index{加速度矢量}是质点 $M$ 的
速度矢量的导数，
$\vect a(t) = \bigl(v_x'(t),\, v_y'(t)\bigr) =
\bigl(x''(t),\, y''(t)\bigr)$，单位为 \unit{m/s^2}：表示速度矢量当前变化的快慢，
以及其变化的方向。
\end{definition}

\begin{remark}[加速而速率不变]\label{rem:g12:kinematics-2d:turning}
$\vect a \neq \vect 0$，只要速度\emph{矢量}发生变化就成立——
大小或\emph{方向}任一变化：刹车的汽车
沿后退方向有加速度；以恒定速率转弯的汽车（沿侧向）同样有加速度。
这正是力所控制的量
（见\cref{ch:g11:forces-and-motion,ch:g12:newtons-laws}）。
\end{remark}

\begin{method}[读取频闪照片]\label{met:g12:kinematics-2d:chrono}
\emph{频闪照片}\index{频闪照片}记录了
运动质点的位置 $M_0, M_1, M_2, \dots$，
时间间隔相等，均为 $\Delta t$。
\begin{enumerate}
  \item 点间距反映速率：间距相等则速率不变；
        间距增大则速率增大；间距减小则速率减小。
  \item $M_i$ 处的速度为对称弦，
        $\vect v_i \approx \vect{M_{i-1}M_{i+1}}/(2\,\Delta t)$,
        画在 $M_i$ 处，与虚线路径相切。
  \item 求 $M_i$ 处的加速度：将 $\vect v_{i-1}$ 与
        $\vect v_{i+1}$ 平移到 $M_i$，画出
        $\Delta\vect v = \vect v_{i+1} - \vect v_{i-1}$ 首尾相接；
        然后 $\vect a_i \approx \Delta\vect v/(2\,\Delta t)$。
\end{enumerate}
\end{method}

\begin{omfigure}
\begin{tikzpicture}[scale=1.05, font=\small]
  \fill (0,2.05) circle (2pt) (0.8,1.95) circle (2pt)
    (1.6,1.72) circle (2pt) (2.4,1.36) circle (2pt)
    (3.2,0.87) circle (2pt) (4.0,0.25) circle (2pt);
  \node[above] at (0,2.05) {$M_0$};
  \node[above] at (1.6,1.72) {$M_2$};
  \node[above right] at (2.4,1.36) {$M_3$};
  \node[above right] at (3.2,0.87) {$M_4$};
  \draw[-Stealth, omDef, thick] (1.6,1.72) -- ++(1.2,-0.44)
    node[right] {$\vect v_2$};
  \draw[-Stealth, omDef, thick] (3.2,0.87) -- ++(1.2,-0.83)
    node[right] {$\vect v_4$};
  \draw[-Stealth, omDef, dashed] (2.4,1.36) -- ++(1.2,-0.44);
  \draw[-Stealth, omDef] (2.4,1.36) -- ++(1.2,-0.83);
  \draw[-Stealth, omThm, very thick] (3.6,0.92) -- ++(0,-0.39)
    node[right] {$\Delta\vect v$};
  \draw[-Stealth, omThm, very thick] (2.4,1.36) -- ++(0,-0.75)
    node[left] {$\vect a$};
\end{tikzpicture}

{\small 抛出小球的频闪照片：$\vect v_2$ 与 $\vect
v_4$ 由弦构造，并平移到 $M_3$；它们的差
$\Delta\vect v$ 给出 $\vect a$——竖直向下。}
\end{omfigure}

\section{直线运动}

沿直线运动时，一条坐标轴即可：$x(t)$，$v(t) = x'(t)$，
$a(t) = v'(t)$，均为代数量。

\begin{definition}[匀速直线运动]\label{def:g12:kinematics-2d:uniform}
称某运动为\emph{匀速直线运动}\index{匀速直线
运动}，当 $\vect v$ 为常矢量时：轨迹为直线，
速率恒定，$\vect a = \vect 0$。于是 $x(t) = v\,t + x_0$。
\end{definition}

\begin{proposition}[匀变速直线运动]\label{prop:g12:kinematics-2d:uarm}
\index{匀变速直线运动}%
若质点沿 $Ox$ 以\emph{恒定}的加速度 $a$ 运动，
从 $t = 0$ 时刻、位置 $x_0$ 处以速度 $v_0$ 出发，则
\[
v(t) = a\,t + v_0,
\qquad
x(t) = \tfrac12\,a\,t^2 + v_0\,t + x_0 .
\]
\end{proposition}

\begin{proof}
$v$ 是常量 $a$ 的一个原函数：$v(t) = a t + C$，且
由 $t = 0$ 得 $C = v_0$。进而 $x$ 是
$a t + v_0$ 的一个原函数：$x(t) = \tfrac12 a t^2 + v_0 t + C'$，其中 $C' = x_0$。
\end{proof}

\begin{omfigure}
\begin{tikzpicture}[font=\small]
\begin{axis}[omaxis, name=ax1, width=8.4cm, height=3.2cm,
    ylabel={$x$ (\unit{m})}, xmin=0, xmax=4, ymin=0, ymax=21]
  \addplot[omDef, thick, domain=0:4, samples=40] {0.5*2*x^2 + x};
\end{axis}
\begin{axis}[omaxis, name=ax2, at={(ax1.below south west)},
    anchor=north west, yshift=-3mm, width=8.4cm, height=3.2cm,
    ylabel={$v$ (\unit{m/s})}, xmin=0, xmax=4, ymin=0, ymax=10]
  \addplot[omThm, thick, domain=0:4, samples=2] {2*x + 1};
\end{axis}
\begin{axis}[omaxis, at={(ax2.below south west)}, anchor=north west,
    yshift=-3mm, width=8.4cm, height=3.2cm, xlabel={$t$ (\unit{s})},
    ylabel={$a$ (\unit{m/s^2})}, xmin=0, xmax=4, ymin=0, ymax=3]
  \addplot[omProp, thick, domain=0:4, samples=2] {2};
\end{axis}
\end{tikzpicture}

{\small 匀变速直线运动（$a = \qty{2.0}{m/s^2}$，
$v_0 = \qty{1.0}{m/s}$，$x_0 = 0$）：每幅图都是
上方那幅图的导数——抛物线、其斜率、斜率的斜率。}
\end{omfigure}

\begin{example}[制动距离]\label{ex:g12:kinematics-2d:braking}
一辆汽车以速率 $v_0$ 行驶，以恒定减速度 $a$ 制动
（加速度为 $-a$）。停车时 $v(t) = -a t + v_0 = 0$，
即 $t_s = v_0/a$ 时，已行驶距离为 $d = -\tfrac12 a t_s^2 + v_0 t_s =
v_0^2/(2a)$。平方关系是道路安全的核心结论：速率加倍，
制动距离\emph{变为四倍}——取 $a =
\qty{6.0}{m/s^2}$ 时，制动距离为 \qty{16}{m}（自 \qty{50}{km/h} 起），
为 \qty{64}{m}（自 \qty{100}{km/h} 起）。
\end{example}

\section{匀速圆周运动与弗勒内标架}

\begin{definition}[匀速圆周运动]\label{def:g12:kinematics-2d:ucm}
称某运动为\emph{匀速圆周运动}\index{匀速圆周运动}，
当其轨迹为半径为 $R$ 的圆且速率 $v$ 恒定时，
绕行一圈所用的时长称为\emph{周期}\index{周期!转动} $T$，
而\emph{频率}\index{频率!
转动} $f = 1/T$（单位为 \unit{Hz}）是每秒绕行的圈数：
$v = 2\pi R/T = 2\pi R f$.
\end{definition}

速率恒定——然而速度在转动：此处\emph{确实}存在
一个加速度。方向如何？大小又是多少？

\begin{theorem}[向心加速度]\label{thm:g12:kinematics-2d:centripetal}
\index{向心加速度}%
在半径为 $R$、速率为 $v$ 的匀速圆周运动中，
这一加速度——称为\emph{向心}加速度——方向由运动质点
指向圆心，其模为常量：
\[
a = \frac{v^2}{R}.
\]
\end{theorem}

\begin{proof}
以圆心为原点建立坐标轴，并令 $\omega = v/R$，于是
弧长 $v t$ 所对圆心角为 $\omega t$，故位置为
$x(t) = R\cos(\omega t)$，$y(t) = R\sin(\omega t)$。求导得
$\vect v = (-R\omega\sin(\omega t),\, R\omega\cos(\omega t))$，其模
恒为 $R\omega = v$；再次求导得
$\vect a = -\omega^2\,\vect{OM}$，方向与位置矢量
相反——指向圆心——其模为 $\omega^2 R = v^2/R$。
\end{proof}

\begin{omfigure}
\begin{tikzpicture}[scale=1.0, font=\small]
  \draw[thick] (0,0) circle (1.7);
  \fill (0,0) circle (1.5pt); \node[below left] at (0,0) {$O$};
  \draw[dashed, omExo] (0,0) -- (215:1.7)
    node[midway, below right] {$R$};
  \fill (35:1.7) circle (2pt); \node[above right] at (35:1.7) {$M$};
  \draw[-Stealth, omDef, thick] (35:1.7) -- ++(-0.69,0.98)
    node[above] {$\vect v$};
  \draw[-Stealth, omThm, very thick] (35:1.7) -- ++(215:0.95)
    node[below right] {$\vect a$};
\end{tikzpicture}

{\small 匀速圆周运动：$\vect v$ 沿切线方向，长度不变；
$\vect a$ 指向圆心，长度恒为 $v^2/R$——始终
只转动速度，而从不拉长它。}
\end{omfigure}

\begin{definition}[弗勒内标架]\label{def:g12:kinematics-2d:frenet}
曲线路径上一点在当前位置的\emph{弗勒内标架}\index{弗勒内标架}
是一对相互垂直的单位矢量
$(\vect{u_t}, \vect{u_n})$：$\vect{u_t}$ 沿切向，沿运动方向；
$\vect{u_n}$ 沿法向，指向曲线该处最贴合的圆的圆心，
该圆半径为 $R$。
\end{definition}

\begin{proposition}[弗勒内标架中的加速度]\label{prop:g12:kinematics-2d:frenetacc}
沿任一平面路径运动，局部半径为 $R$，速率为 $v(t)$，
\[
\vect a = a_t\,\vect{u_t} + a_n\,\vect{u_n},
\qquad a_t = \frac{\dd v}{\dd t},
\qquad a_n = \frac{v^2}{R}:
\]
\emph{切向加速度}\index{切向加速度}
$a_t$ 改变速率的大小；\emph{法向加速度}\index{法向加速度}
$a_n$ 改变速度的方向。
\end{proposition}
\admitted

\begin{remark}[证明的出处]\label{rem:g12:kinematics-2d:frenetproof}
已证明两种极端情形：直线上的纯 $a_t$
（\cref{prop:g12:kinematics-2d:uarm}），圆周上的纯 $a_n$
（\cref{thm:g12:kinematics-2d:centripetal}）；一般分解
在本科一年级分册中推导。
\end{remark}

\begin{example}[刹车进入弯道]\label{ex:g12:kinematics-2d:bend}
一辆汽车驶入半径为 $R = \qty{50}{m}$ 的弯道，速率为 $v = \qty{15}{m/s}$
并同时刹车，$\dd v/\dd t = \qty{-3.0}{m/s^2}$：$a_t =
\qty{-3.0}{m/s^2}$，$a_n = 15^2/50 = \qty{4.5}{m/s^2}$，于是 $a =
\sqrt{3.0^2 + 4.5^2} \approx \qty{5.4}{m/s^2}$，方向既向后
\emph{又}指向弯道内侧。轮胎必须同时提供这两者——
这正是赛车手要\emph{在入弯前}刹车的原因。
\end{example}

\section{练习}

\begin{exercise}[$\star$]\label{exo:g12:kinematics-2d:1}
已知 $x(t) = 3.0\,t$，$y(t) = 4.0\,t$（米、秒），
求 $\vect v$、速率、轨迹以及运动类型。
\end{exercise}

\begin{exercise}[$\star$]\label{exo:g12:kinematics-2d:2}
在每隔 \qty{0.10}{s} 拍摄一次的频闪照片上，各点恰好落在一条
直线上，相邻两点间距均为 \qty{0.35}{m}。判断该运动类型，并
计算速率。若间距逐渐增大，说明什么？
\end{exercise}

\begin{exercise}[$\star$]\label{exo:g12:kinematics-2d:3}
已知 $x(t) = 2.0\,t^2$ 沿 $Ox$ 方向运动（米、秒），求 $v(t)$
和 $a(t)$，指出运动类别，并计算
$t = \qty{3.0}{s}$ 时的速率。
\end{exercise}

\begin{exercise}[$\star$]\label{exo:g12:kinematics-2d:4}
一匹旋转木马上的木马沿半径为 \qty{4.5}{m} 的圆周运动，每
\qty{8.0}{s} 转一圈：求频率、速率和加速度
（大小与方向）。
\end{exercise}

\begin{exercise}[$\star$]\label{exo:g12:kinematics-2d:5}
$v(t)$ 图像呈直线上升，从 $t = 0$ 时的 \qty{2.0}{m/s}
升至 $t = \qty{4.0}{s}$ 时的 \qty{10.0}{m/s}。读出加速度，
再求所走路程（图线下的面积）。
\end{exercise}

\begin{exercise}[$\star\star$]\label{exo:g12:kinematics-2d:6}
一列高速列车通过半径为 \qty{6.0}{km} 的弯道，
速率为 \qty{320}{km/h}。求 $a_n$，并与 $g$ 比较；
即使弯道设计得完全合适，高速铁路为何仍须避开急弯？
\end{exercise}

\begin{exercise}[$\star\star$]\label{exo:g12:kinematics-2d:7}
一辆小车沿直轨道运动，其 $x(t) = -1.2\,t^2 + 8.0\,t + 3.0$
（米、秒）。求 $v(t)$ 和 $a$，以及换向的时刻与位置，
并描述换向前后的运动。
\end{exercise}

\begin{exercise}[$\star\star$]\label{exo:g12:kinematics-2d:8}
沿一条直线拍摄的频闪照片，每隔 $\Delta t = \qty{0.20}{s}$ 取一个位置，
给出 $x = 0$，$0.10$，$0.28$，$0.54$，$\qty{0.88}{m}$。估算
中间三个点处的速度，
（见\cref{met:g12:kinematics-2d:chrono}），再求加速度；作出结论。
\end{exercise}

\begin{exercise}[$\star\star$]\label{exo:g12:kinematics-2d:9}
月球绕地球近似做圆周运动，轨道半径为 \qty{3.84e8}{m}，
周期为 \qty{27.3}{d}。求其速率和向心加速度；
将后者与 $g$ 比较，已知月球距地心约 $60$ 个地球半径
远（见\cref{ch:g10:universal-gravitation}）。
\end{exercise}

\begin{exercise}[$\star\star$]\label{exo:g12:kinematics-2d:10}
一台洗衣机滚筒，半径为 \qty{25}{cm}，以 $1200$
转/分的转速旋转。求 $f$、$T$、边缘速率以及边缘
加速度（以 $g$ 的倍数表示）。为什么衣物会贴在筒壁上，
而水却被甩出？
\end{exercise}

\begin{exercise}[$\star\star$]\label{exo:g12:kinematics-2d:11}
一辆汽车以 \qty{90}{km/h} 的速率行驶，以恒定加速度 $a = \qty{7.5}{m/s^2}$ 刹车：
求停车时间和制动距离，并画出 $v(t)$ 的图像。
\end{exercise}

\begin{exercise}[$\star\star\star$]\label{exo:g12:kinematics-2d:12}
对于 $x(t) = R\cos(\omega t)$，$y(t) = R\sin(\omega t)$，求
$\vect v$；验证速率为常量 $R\omega$，并验证
$\vect v \perp \vect{OM}$；求 $\vect a$，证明它为
$-\omega^2\,\vect{OM}$，并由此推出 $a = v^2/R$。
\end{exercise}

\begin{exercise}[$\star\star\star$]\label{exo:g12:kinematics-2d:13}
一辆摩托车进入半径为 \qty{80}{m} 的弯道，速率为 \qty{20}{m/s}，
并以 \qty{2.5}{m/s^2} 的加速度制动。求 $a_t$、$a_n$，以及
$\vect a$ 的模和它与运动方向的夹角。轮胎的抓地力
总计至多为 \qty{7.0}{m/s^2}，骑手是否在极限之内？
\end{exercise}

\begin{exercise}[$\star\star\star$]\label{exo:g12:kinematics-2d:14}
一列地铁从静止出发，以 \qty{1.2}{m/s^2} 的加速度加速，直至速率达到 \qty{20}{m/s}，
匀速行驶 \qty{40}{s}，随后以 \qty{1.5}{m/s^2} 的加速度制动至停止。
求各阶段的时长和距离，以及总距离与
平均速率。画出 $v(t)$ 的图像，并用面积检验距离。
\end{exercise}

\begin{exercise}[$\star\star\star$]\label{exo:g12:kinematics-2d:15}
一颗地球静止卫星绕地心做圆周运动，轨道半径为 \qty{4.22e7}{m}，
周期为 \qty{86164}{s}。求其速率和
向心加速度。是何物提供了它？卫星为何
能始终位于赤道上空某一固定点正上方（见\cref{ch:g12:satellites-kepler}）？
\end{exercise}

\section{问题：事故报告}

\begin{problem}[{周末问题——事故报告：逐帧读取行车记录仪影片，
由刹车痕迹重建制动规律，
由环岛作证抓地力，并以千米每小时为单位给出
结论}]
\label{pb:g12:kinematics-2d:1}
一辆汽车撞上了一辆从支路驶出的面包车；无人受伤，
但保险公司各执一词。调查人员掌握了高架摄像头
拍摄的录像（每隔 \qty{0.50}{s} 一帧）、实测的刹车痕迹，以及
该车的行车记录仪。限速为 \qty{50}{km/h}。

\medskip
\noindent\textbf{第一部分——读取影片。}
反应发生之前的各帧沿直线道路给出以下数据：
\begin{center}\small
\begin{tabular}{lccccc}
$t$ (\unit{s}) & 0.00 & 0.50 & 1.00 & 1.50 & 2.00\\
$x$ (\unit{m}) & 0.0 & 9.0 & 18.0 & 27.0 & 36.0
\end{tabular}
\end{center}
\begin{enumerate}
  \item 为何报告必须首先指明参考系和坐标轴？
        本题采用的是哪一组参考系与坐标轴？
  \item 求四个时间段内各自的平均速度。
  \item 这是哪一类运动？请根据表中数据说明理由。
  \item 求汽车的速率 $v_0$，分别以 \unit{m/s} 和 \unit{km/h} 表示。
  \item 这两秒内的加速度是多少？
\end{enumerate}

\medskip
\noindent\textbf{第二部分——刹车痕迹。}
此干燥沥青路面上的轮胎测试给出恒定的制动加速度大小，
$a = \qty{6.0}{m/s^2}$；刹车痕迹长为 $d_b = \qty{27}{m}$。
取制动开始时为 $t = 0$、$x = 0$，速率 $v_B$ 未知。
\begin{enumerate}[resume]
  \item 通过求原函数，写出制动过程中的 $v(t)$ 和 $x(t)$。
  \item 求停车时间 $t_s$；证明 $d_b = v_B^2/(2a)$。
  \item 由刹车痕迹推出 $v_B$。它与第一部分一致吗？
  \item 反应时间为 \qty{1.0}{s}，求反应距离和
        总停车距离。
  \item 由第 7 题说明，速率加倍为何使
        制动距离变为四倍；求恰为 \qty{50}{km/h} 时的 $d_b$。
\end{enumerate}

\medskip
\noindent\textbf{第三部分——环岛作证。}
此前，行车记录仪显示汽车以恒定速率绕行一条环岛车道，
半径为 $R = \qty{12.5}{m}$，转过四分之一圈
用时 \qty{2.5}{s}。路面至多能提供 $a_n = \mu g$，
其中 $\mu = 0.70$。
\begin{enumerate}[resume]
  \item 此处速率恒定——速度是否也恒定？汽车
        是否在加速？请说明理由。
  \item 求汽车在环岛上的速率。
  \item 求其向心加速度（大小和方向）。
  \item 求不打滑的最大速率。汽车是否在此速率之内？
  \item 求以该速率绕行一周的周期和频率。
\end{enumerate}

\medskip
\noindent\textbf{第四部分——结论。}
驾驶员开始反应的瞬间，面包车出现在汽车前方
$D = \qty{41}{m}$ 处。
\begin{enumerate}[resume]
  \item 制动前的速率超出限速多少（分别以 \unit{km/h} 和
        百分比表示）？
  \item 求撞击时汽车的速率。
  \item 在恰为 \qty{50}{km/h} 时重做停车计算（反应
        和制动条件不变）：汽车能否停下？余量为多少？
  \item 给出一般的停车距离公式 $d(v)$；为何
        速率在其中出现两次——一次是线性的，一次是平方的？
  \item 结论：用一句话，并附上有关数据，说明驾驶员
        当时的速率是多少？若以限速行驶，事故是否会发生？
\end{enumerate}
\end{problem}
